\documentclass[11pt]{article}
\usepackage[margin=1in]{geometry}
\usepackage{amsmath,amssymb,amsthm}
\usepackage{hyperref}
\usepackage{xurl}

\newtheorem{theorem}{Theorem}
\newtheorem{lemma}[theorem]{Lemma}
\theoremstyle{definition}
\newtheorem{definition}[theorem]{Definition}
\theoremstyle{remark}
\newtheorem{remark}[theorem]{Remark}

\let\oldus\_
\renewcommand{\_}{\oldus\allowbreak}

\title{A disproof of Conjecture 00000000172\\[2pt]
\large Under the stated definition, $\sigma = \lim \sigma_n^{1/n} = 1$}
\date{}

\begin{document}
\sloppy
\maketitle

\section{The conjecture and its reading}

Conjecture 00000000172 reads: \emph{Definition: Somos's constant $\sigma = \lim \sigma_n^{1/n}$, where
$\sigma_n = \sigma_{n-1} + \sigma_{\lfloor n/2\rfloor}$, $\sigma_1 = 1$. Conjecture: The density of 1's in the
binary expansion of $\sigma$ diverges, with counting function growing as $\Theta(\log\log n)$; this
yields an explicit proof path for the irrationality of $\sigma$.} The Chinese text gives the same
definition ($\sigma=\lim \sigma_n^{1/n}$, $\sigma_n=\sigma_{n-1}+\sigma_{\lfloor n/2\rfloor}$, $\sigma_1=1$) and the
same three claims.

\paragraph{Which $\sigma$.} We use the definition written in the conjecture (identical in both
languages). The recurrence is OEIS A033485, ``$a(n) = a(n-1) + a(\lfloor n/2\rfloor)$, $a(1) = 1$''
(retrieved), with terms $1,2,3,5,7,10,13,18,\dots$; it is defined for $n\ge1$, the recurrence applying
for $n \ge 2$. We note openly that the name ``Somos' constant'' usually denotes a different number,
Somos' quadratic recurrence constant $\sqrt{1\sqrt{2\sqrt{3\cdots}}} \approx 1.661687949633594$
(Wikipedia, retrieved); nothing below concerns that number.

\paragraph{The claims refuted.} Writing $C(n)$ for the number of 1's among the first $n$ binary digits of
$\sigma$:
(i) $C(n) = \Theta(\log\log n)$;
(ii) the density $C(n)/n$ diverges (we refute both meanings: ``has no limit'' and ``tends to $+\infty$'');
(iii) $\sigma$ is irrational (the claim that the above ``yields a proof'' of irrationality).
We show $\sigma = 1$; then (iii) fails, and (i), (ii) fail for every binary expansion of $1$.

\paragraph{Conventions.} $\log$ is the natural logarithm; $\sigma_n^{1/n}$ is the real power of the
positive number $\sigma_n$. A binary expansion of $x$ is a pair $(m,d)$ with $m\in\mathbb Z$, digits
$d_k \in\{0,1\}$ ($k\ge0$) and $x = m + \sum_{k\ge0} d_k 2^{-(k+1)}$. We set
$C_d(n) = \#\{k<n : d_k = 1\}$ (fractional digits). To cover any convention for counting the digits of
the integer part, every statement is proved for $c + C_d(n)$ with an arbitrary real offset $c$.
$\Theta$ is Mathlib's \texttt{IsTheta} (both $f = O(g)$ and $g = O(f)$ along \texttt{atTop}).

\section{Proof that $\sigma = 1$}

Let $a:\mathbb N\to\mathbb R$ be \emph{any} sequence with $a_1 = 1$ and $a_n = a_{n-1} + a_{\lfloor n/2\rfloor}$
for $n\ge 2$ (\texttt{IsSomosSeq}; $a_0$ is unconstrained and never used). Such a sequence exists
(\texttt{A}, \texttt{isSomosSeq\_A}).

\begin{lemma}\label{l:basic}
(\texttt{one\_le}) $a_n \ge 1$ for $n\ge1$. (\texttt{mono}) $a_m \le a_n$ for $1\le m\le n$.
(\texttt{le\_mul\_half}) $a_n \le n\,a_{\lfloor n/2\rfloor}$ for $n\ge 2$.
(\texttt{le\_pow}) $a_n \le n^{\lfloor \log_2 n\rfloor + 1}$ for $n \ge 1$.
\end{lemma}
\begin{proof}
The first two are inductions using $a_{\lfloor n/2\rfloor} \ge 1$. For the third, $a_2 = 2 = 2a_1$, and if
$a_n \le n a_{\lfloor n/2\rfloor}$ then
$a_{n+1} = a_n + a_{\lfloor (n+1)/2\rfloor} \le n a_{\lfloor (n+1)/2\rfloor} + a_{\lfloor (n+1)/2\rfloor}$ by
monotonicity. For the fourth, by strong induction, using
$\lfloor\log_2\lfloor n/2\rfloor\rfloor + 1 = \lfloor \log_2 n\rfloor$ for $n\ge 2$ (Mathlib
\texttt{Nat.log\_div\_base}): $a_n \le n\cdot \lfloor n/2\rfloor^{\lfloor\log_2 n\rfloor}\le n^{\lfloor\log_2 n\rfloor+1}$.
\end{proof}

\begin{theorem}[\texttt{tendsto\_root}, \texttt{sigma\_eq\_one}]
$a_n^{1/n} \to 1$. Hence the conjecture's $\sigma$ exists and equals $1$.
\end{theorem}
\begin{proof}
Since $2^{\lfloor\log_2 n\rfloor}\le n$, $\lfloor\log_2 n\rfloor \le \log n/\log 2$, so by Lemma~\ref{l:basic}
(\texttt{log\_le}) $0 \le \log a_n \le (\log n/\log 2 + 1)\log n$. Thus
$0 \le \frac{\log a_n}{n} \le \frac{(\log n)^2}{(\log 2)\, n} + \frac{\log n}{n} \to 0$ (Mathlib
\texttt{Real.tendsto\_pow\_log\_div\_mul\_add\_atTop}), and $a_n^{1/n} = \exp(\log a_n / n) \to e^0 = 1$.
\end{proof}

\section{Failure of the claims}

\begin{lemma}[\texttt{expansion\_of\_one}]
If $(m,d)$ is a binary expansion of $1$, then either $d_k = 0$ for all $k$, or $d_k = 1$ for all $k$.
\end{lemma}
\begin{proof}
Let $S = \sum_k d_k 2^{-(k+1)}$; the terms are non-negative and bounded by $2^{-(k+1)}$, so
$0 \le S\le 1$ and $m = 1 - S \in [0,1]\cap\mathbb Z$. If $m = 1$ then $S = 0$, and a digit $d_k = 1$
would give $S \ge 2^{-(k+1)} > 0$. If $m = 0$ then $S = 1$, so
$\sum_k (2^{-(k+1)} - d_k 2^{-(k+1)}) = 1 - 1 = 0$ with non-negative terms, and a digit $d_k = 0$ would
make this sum $\ge 2^{-(k+1)} > 0$.
\end{proof}

So $C_d(n) = 0$ for all $n$, or $C_d(n) = n$ for all $n$.

\begin{lemma}[\texttt{refute}]
Let $c\in\mathbb R$ and $F(n) = c$ for all $n$, or $F(n) = c+n$ for all $n$. Then $F$ is not
$\Theta(\log\log n)$, and $F(n)/n$ converges.
\end{lemma}
\begin{proof}
If $F \equiv c$ and $\log\log n = O(c)$, then $\log\log n$ would be eventually bounded, but
$\log\log n\to\infty$. If $F(n) = c+n$ and $c + n = O(\log\log n)$, then, as
$\log\log n = o(\log n)$ and $\log n = o(n)$ (Mathlib \texttt{Real.isLittleO\_log\_id\_atTop}),
$c + n = o(n)$, so eventually $|c+n|\le n/2$, false for $n > 2|c|$. Finally $c/n\to0$ and
$(c+n)/n = 1 + c/n\to 1$.
\end{proof}

\begin{theorem}[\texttt{C172.conjecture\_172\_false}]
Let $a$ satisfy the definition and let $a_n^{1/n}\to\sigma$. Then $\sigma = 1$; $\sigma$ is not irrational;
and for every binary expansion $(m,d)$ of $\sigma$ and every $c \in \mathbb R$: $c + C_d(n)$ is not
$\Theta(\log\log n)$; $(c + C_d(n))/n$ converges to a finite limit; and $(c + C_d(n))/n\not\to+\infty$.
\end{theorem}
\begin{proof}
Uniqueness of limits and the theorem above give $\sigma=1$; $1$ is rational (Mathlib
\texttt{not\_irrational\_one}); the two lemmas give the rest.
\end{proof}

\begin{remark}[Scope]
All three claims (i)--(iii) are refuted for every binary expansion of $\sigma$. A reading of ``the density
diverges'' as ``the number of 1's tends to infinity'' (rather than the density) is not refuted on its own
for the expansion $0.111\ldots$; it is refuted together with clause (i), which fails for every expansion.
\end{remark}

\section{Lean formalization}

File \texttt{lean/Conjecture172/Basic.lean} (256 lines, only \texttt{import Mathlib}), namespace
\texttt{C172}. Main theorem (ASCII rendering of the Unicode source):
\begin{verbatim}
theorem conjecture_172_false (a : Nat -> Real) (ha : IsSomosSeq a)
    (s : Real)
    (hs : Tendsto (fun n : Nat => a n ^ ((1 : Real) / n)) atTop (nhds s)) :
    s = 1 /\ not (Irrational s) /\
    forall (m : Int) (d : Nat -> Nat), IsBinaryExpansion s m d ->
    forall c : Real,
      not ((fun n : Nat => c + (onesCount d n : Real)) =Theta[atTop]
            fun n : Nat => log (log (n : Real))) /\
      (exists L, Tendsto (fun n : Nat => (c + (onesCount d n : Real)) / n)
                   atTop (nhds L)) /\
      not (Tendsto (fun n : Nat => (c + (onesCount d n : Real)) / n)
                   atTop atTop)
\end{verbatim}
Other theorems: \texttt{isSomosSeq\_A}, \texttt{one\_le}, \texttt{mono}, \texttt{le\_mul\_half},
\texttt{le\_pow}, \texttt{log\_le}, \texttt{tendsto\_root}, \texttt{expansion\_of\_one},
\texttt{tendsto\_loglog}, \texttt{loglog\_isLittleO}, \texttt{refute}, \texttt{sigma\_eq\_one}.

\paragraph{Axiom audit.} \texttt{lake build} succeeds, and \texttt{\#print axioms} for
\texttt{C172.conjecture\_172\_false}, \texttt{C172.sigma\_eq\_one} and \texttt{C172.isSomosSeq\_A}
reports only \texttt{[propext, Classical.choice, Quot.sound]}. No \texttt{sorry}, no
\texttt{native\_decide}, no added axiom.

\section*{Sources}
\begin{itemize}
\item OEIS A033485, \url{https://oeis.org/A033485}: ``A033485: a(n) = a(n-1) + a(floor(n/2)), a(1) = 1.''
\item Wikipedia, \emph{Somos' quadratic recurrence constant},
\url{https://en.wikipedia.org/wiki/Somos\%27_quadratic_recurrence_constant}: ``which gives a numerical value
of approximately'' $\sigma = 1.661687949633594121295\dots$
\end{itemize}

\end{document}
